% Section 5.6 -- structural breadth, one paragraph. Direction in the main text,
% magnitudes and masses in Appendix E. Reported as unstable, NOT as evidence.

\subsection{Structural breadth}
\label{sec:breadth}

Breadth, as defined in Section~\ref{sec:definitions}, counts the distinct
communities of the partition a proof's premises touch, per tactic step. Machine
proofs touch fewer communities per step than human ones at every scope we can
measure, and the direction is the only part of that we report: the gap is
\breadthGapFull{} on the full corpora, \breadthGapSupport{} under common
support and \breadthGapMatched{} matched 1:1, the last two computed on
\supportRetentionMachine{} of the machine corpus. The magnitude therefore moves
by a factor of four across designs that move the coverage result in its fourth
decimal place, and an earlier sign flip was an artifact of a region cut on
character length, which Section~\ref{sec:length-axis} rejects as a formatting
control. We report the metric as unstable, having defined it before we looked,
and we do not offer it as evidence for or against the coverage result, which
stands on its own controls. Appendix~\ref{app:breadth} carries the masses, the
excluded region and the partition skew that keep it secondary.
